\section{Spin one in both readings}
\label{sec:spin}

A photon is a massless particle of spin one with two helicities. Its covariant description is a field strength, a two-form, and
its potential is a four-vector defined up to gauge. The arena offers two places for it. In the \emph{interior reading} the
potential is a grade-one element of the four-plane $E$ on which the Dirac sector is built and the field strength is the grade-two
module of \cref{prop:bivectorF}. In the \emph{frame reading} the field strength occupies the six directions themselves, under the
algebra $\lG$ of \cref{thm:GLorentz}. This section shows that each reading produces a free field with the spin content of the photon,
that the two are related by the Klein correspondence, and that at the second rung both can be housed under one Lorentz action and
coupled to a Dirac block. \Cref{fig:tworoutes} summarises the comparison.

\subsection{The univectorial photon}

\begin{proposition}[Grade one under the interior action]
\label{prop:gradeone}
Under $\lint\oplus\uu1_J$, with $J=e_2e_3$ acting by $\ad(J/2)$,
\[
V=\bigl(\tfrac12,\tfrac12\bigr)_0\oplus(0,0)_{+1}\oplus(0,0)_{-1}.
\]
The Lorentz Casimir is $\operatorname{diag}(3,0,0,3,3,3)$ on $(e_1,\dots,e_6)$, and $\ad(J/2)$ vanishes on the four-plane and acts on
$(e_2,e_3)$ by $\bigl(\begin{smallmatrix}0&1\\-1&0\end{smallmatrix}\bigr)$, with eigenvalues $\pm i$.
\end{proposition}

The potential $A=A^\mu\gamma_\mu$ is therefore neutral under the charge, as a photon must be, while the surplus pair is a Lorentz
scalar of charge $\pm1$. The photon potential is the four-plane part of grade one, not the whole of it.

\begin{proposition}[Spin one and spin one-half from the same rotors]
\label{prop:samerotors}
The interior rotations $S_k$ satisfy $[S_i,S_j]=\epsilon_{ijk}S_k$. On grade one $\ad S_k$ acts on $(e_4,e_5,e_6)$ as the standard
rotation generator $r_k$ and annihilates $e_1,e_2,e_3$, so that $-\sum_k(\ad S_k)^2=\operatorname{diag}(0,0,0,2,2,2)$ and
$\spec\ad S_3=\{\pm i,0,0,0,0\}$. On $S=\R^8$, by left multiplication, $-\sum_kS_k^2=\tfrac34\id$ and $S_3$ has eigenvalues $\pm\tfrac i2$.
Consequently $\exp(2\pi\ad S_3)=\id$ on grade one while $\exp(2\pi S_3)=-\id$ on $S$.
\end{proposition}

A vector transforms by conjugation, $A\mapsto RAR^{-1}$ with $R=e^{\theta S_3}$, and a spinor by left multiplication,
$\psi\mapsto R\psi$. At $\theta=2\pi$ the rotor is $-\id$; the sign cancels in the first case and survives in the second. This is the
space-time-algebra account of the difference between vectors and spinors \cite{Hestenes66}, realised here by the rotors of the Dirac
sector: the same $S_k$ give spin one to the potential and spin one-half to matter.

\begin{theorem}[Maxwell's equations as one Clifford equation]
\label{thm:nablaF}
Put $\nabla=\gamma^\mu\partial_\mu$ and $\sigma_k=\gamma_k\gamma_0$.
\begin{enumerate}[label=(\roman*),leftmargin=2.2em]
\item For $A=A^\mu\gamma_\mu$, $\nabla A=\partial\cdot A+\nabla\wedge A$, and for $F:=\nabla\wedge A$ the grade-one part of $\nabla F$ is
$(\Box A^\nu-\partial^\nu\partial\cdot A)\gamma_\nu$, while its grade-three part vanishes identically.
\item $F=\bE+\Omega\,\bB$ with $\bE=\sum_kE_k\sigma_k$, $E_k=-\partial_kA^0-\partial_tA^k$ and $\bB=\curl\bA$.
\item For every $F=\sum_k(E_k\sigma_k+B_k\Omega\sigma_k)$,
\[
\nabla F=(\nabla\cdot\bE)\gamma_0+(\curl\bB-\partial_t\bE)^k\gamma_k+\Omega\bigl[-(\nabla\cdot\bB)\gamma_0+(\partial_t\bB+\curl\bE)^k\gamma_k\bigr],
\]
with no other grades.
\end{enumerate}
Hence $\nabla F=\mathcal J$, with $\mathcal J=\rho\gamma_0+\mathcal J^k\gamma_k$, is Gauss's law, the Amp\`ere--Maxwell law, the absence of
magnetic charge and Faraday's law, and for $F=\nabla\wedge A$ the last two hold identically.
\end{theorem}

This is Hestenes' single Maxwell equation \cite{Hestenes66}, here on the four-plane of the Dirac sector inside the $8\times8$ matrices
of $\Cl(3,3)$. Plane waves show how the gauge freedom removes the unphysical polarisations and how helicity appears.

\begin{example}[A plane wave of the univectorial photon]
\label{ex:univectorwave}
Let $k=\gamma_0+\gamma_3$, so $k^2=0$, and $\varepsilon=\varepsilon^\mu\gamma_\mu$ with $k\cdot\varepsilon=0$, that is $\varepsilon^0=\varepsilon^3$. Then
\[
F(\varepsilon)=k\varepsilon=\varepsilon^1F_1+\varepsilon^2F_2,\qquad F_j=k\gamma_j ,
\]
because the component of $\varepsilon$ along $k$ drops out, $k^2=0$. The condition $k\cdot\varepsilon=0$ has removed the component along the
other null direction $\gamma_0-\gamma_3$, so only the two transverse polarisations remain, and the gauge freedom
$\varepsilon\mapsto\varepsilon+ck$ is exactly the helicity-zero direction. Moreover $F_j^2=0$ and $F_1F_2+F_2F_1=0$, so every $F(\varepsilon)$ squares
to zero: in the space-time algebra $F^2=\bE^2-\bB^2+2\Omega\,\bE\cdot\bB$, so this is the null condition $|\bE|=|\bB|$,
$\bE\perp\bB$. On $\operatorname{span}\{F_1,F_2\}$,
\[
\ad S_3=\begin{pmatrix}0&-1\\1&0\end{pmatrix},\qquad \Omega=\begin{pmatrix}0&1\\-1&0\end{pmatrix}=-\ad S_3 ,
\]
both with eigenvalues $\pm i$. For positive-frequency waves the eigenvectors of $\ad S_3$ with eigenvalues $\mp i$ carry helicity
$\pm1$, and on them the Hodge star is $i$ times the helicity: circular polarisation is self-duality or anti-self-duality of $F$.
\end{example}

\begin{proposition}[Every field-strength module under the interior action]
\label{prop:censusint}
The Lorentz Casimir on the grades of $\Cl(3,3)$, the number of field-strength modules (Casimir value $4$) and their charges under
$\ad(J/2)$ are
\begin{center}
\begin{tabular}{@{}clcc@{}}
\toprule
grade & spectrum of the Casimir (multiplicity) & field-strength modules & charges\\
\midrule
$0,6$ & $0\,(1)$ & $0$ & ---\\
$1,5$ & $3\,(4)$, $0\,(2)$ & $0$ & ---\\
$2,4$ & $4\,(6)$, $3\,(8)$, $0\,(1)$ & $1$ & $0$\\
$3$ & $4\,(12)$, $3\,(8)$ & $2$ & $\pm1$\\
\bottomrule
\end{tabular}
\end{center}
and on $S$ the Casimir is $\tfrac32$ throughout.
\end{proposition}

Because $J$ commutes with the four-plane, the grade-four module is $J$ times the grade-two one, and $J$ is the complex structure at
one constituent (\cref{thm:i}). Up to that complex structure the neutral field strength is unique and sits at grade two; the two
grade-three copies carry charge $\pm1$ and are not candidates for the photon.

\subsection{The frame photon}

In the frame reading the field strength is the vector $X$ itself, and the boosts of $\lG$ turn out to be the spin matrices of the
photon.

\begin{proposition}[Spin matrices on the frame]
\label{prop:spinmatrices}
$K_k=GL_k$ for $k=1,2,3$, both families commute with $G$ and preserve $\eta$, and the Casimir $\sum K_k^2-\sum L_k^2$ equals $4\,\id$ on
$V$, the field-strength value. On $z=\bB+i\bE$ the rotation $L_k$ acts as $r_k$ and the boost $K_k$ as
\[
s_k=i\,r_k,\qquad (s_k)_{ab}=-i\epsilon_{kab},
\]
which are Hermitian with $[s_i,s_j]=i\epsilon_{ijk}s_k$, $\sum_ks_k^2=2\,\id$ and $\spec s_3=\{-1,0,1\}$. With $p=-i\nabla$ one has
$(s\cdot p)z=\curl z$ identically, so the frame Maxwell system \eqref{eq:RS} is the photon wave equation
\[
i\partial_tz=(s\cdot p)\,z,\qquad p\cdot z=0 ,
\]
and the kernel of $s_3$ is spanned by $\hat e_3$: transversality removes exactly the helicity-zero component of a wave moving along
$\hat e_3$.
\end{proposition}

These are the spin-one matrices that Oppenheimer and Bialynicki-Birula use to write the photon wave function
\cite{Oppenheimer1931,BialynickiBirula96}. The rotations act on $z$ as $r_k$, and their Hermitian generators $s_k=ir_k$ are the action
of the boosts, because on this module a boost is $G$ times a rotation.

\begin{proposition}[Orthogonality on the frame]
\label{prop:orthframe}
For $X=\sum_i(B_ie_i+E_ie_{i+3})$,
\[
z\cdot z=\eta(X,X)-i\,\eta(X,GX).
\]
Hence $\bE\perp\bB$ exactly when $X\perp_\eta GX$, $|\bE|=|\bB|$ exactly when $X$ is $\eta$-null, and both hold exactly when the plane
$\operatorname{span}\{X,GX\}$ is totally $\eta$-null.
\end{proposition}

\begin{theorem}[The axis of a null field]
\label{thm:axis}
For real $a,b\in\R^3$,
\[
\bigl(s\cdot(a\times b)\bigr)(a+ib)=|b|^2a+i|a|^2b-(a\cdot b)(b+ia).
\]
Consequently a nonzero $z=\bB+i\bE$ satisfies $z\cdot z=0$ if and only if $(s\cdot n)z=z$ for some real unit vector $n$, and then
$n=\bB\times\bE/|\bB\times\bE|$ uniquely.
\end{theorem}

\begin{proof}
If $z\cdot z=0$ then $\bB\cdot\bE=0$ and $|\bB|=|\bE|\neq0$; the identity gives $(s\cdot(\bB\times\bE))z=|\bB|^2z$ and
$|\bB\times\bE|=|\bB|^2$. Conversely $n\cdot((s\cdot n)z)=0$ and $z\cdot((s\cdot n)z)=0$ identically, and the characteristic polynomial
of $s\cdot n$ is $-\lambda(\lambda^2-|n|^2)$, so the eigenvalue $1$ forces $n\cdot z=0$ and $z\cdot z=0$. For uniqueness,
$(s\cdot n)z=in\times z=z$ splits into $\bB=-n\times\bE$ and $\bE=n\times\bB$, whence $n\perp\bB$ and $\bB\times\bE=|\bB|^2n$.
\end{proof}

Since $\bB\times\bE=\bE\times\Btext$ is the Poynting vector, a null field is the $m=+1$ state of the photon spin along its own energy
flux, the frame form of the classical fact that a null field determines a unique direction of propagation
\cite{Synge56,Robinson61}. Orthogonality with equal magnitudes therefore singles out the null fields and their axis. It does not
single out a helicity.

\begin{example}[Three null waves with one axis]
\label{ex:threewaves}
The fields
\[
z_\pm=(1,i,0)\,e^{\pm ik(x_3-t)},\qquad z_{\rm lin}=(1,i,0)\cos k(x_3-t)
\]
all solve \eqref{eq:RS}. In each, $\bB\cdot\bE=0$, $|\bB|=|\bE|$ and $\bB\times\bE$ points along $+\hat e_3$ wherever $z\neq0$, so each is an
$m=+1$ state along $\hat e_3$. At $x_3=0$ the electric field turns according to $(\bE\times\dot\bE)\cdot\hat e_3=+k,\,-k,\,0$: helicity
$+1$, helicity $-1$, and a linear superposition. The difference is the sign of the frequency, $z_+\propto e^{-ikt}$ against
$z_-\propto e^{+ikt}$, which is the Bialynicki-Birula assignment for the Riemann--Silberstein vector \cite{BialynickiBirula96}. Not every
radiation field is null either: $z=(1,i,0)\cos k(x_3-t)+(0,1,i)\cos k(x_1-t)$ solves the system and has $|\bB|=|\bE|$, but
$\bB\cdot\bE=\cos k(x_1-t)\cos k(x_3-t)$.
\end{example}

Orthogonality alone is weaker still. The field $X=(\bB;\bE)=(1,0,0;\,0,2,0)$ has $\eta(X,GX)=0$, so $\bE\perp\bB$, but
$\eta(X,X)=1-4=-3$; the boost $e^{\theta K_3}$ with $\theta=-\tfrac12\log3$ carries it to $(0,0,0;\,0,\sqrt3,0)$, a purely electric field
with the same invariant. A field with $\bE\perp\bB$ and $|\bE|\neq|\bB|$ is purely electric or purely magnetic in a suitable frame
\cite{LandauLifshitz2}. On the null line $\C(1,i,0)$, by contrast, $e^{\theta K_3}$ acts by multiplication by $e^\theta$; with
$\theta=-\operatorname{artanh}\beta$ for a frame moving with velocity $\beta\hat e_3$ this is the Doppler factor $\sqrt{(1-\beta)/(1+\beta)}$.
The stabiliser of the axis acts on a null field by the helicity phase and by the Doppler factor.

\subsection{Potentials on the frame, and helicity as a charge}
\label{sec:potentials}

Place a vector potential on each triple, $\bA$ on $e_1,e_2,e_3$ and $\mathbf C$ on $e_4,e_5,e_6$, write $X_A=(\bA;\mathbf C)$ and
$X_F=\curl X_A$, the curl acting on each triple, and take
\begin{equation}
\label{eq:SSL}
\mathcal L=\tfrac12\,(\curl X_A)\cdot G\dot X_A-\tfrac12\,|\curl X_A|^2
\end{equation}
with the Euclidean dot product of $\R^6$, for which $G$ is orthogonal.

\begin{theorem}[Duality-symmetric Maxwell theory on the frame]
\label{thm:SS}
\begin{enumerate}[label=(\roman*),leftmargin=2.2em]
\item The Euler--Lagrange expression of \eqref{eq:SSL} is $G\curl\bigl(\dot X_A+G\curl X_A\bigr)$. Every solution of the first-order
equation $\dot X_A=-G\curl X_A$ solves it, and on a contractible region every solution satisfies the first-order equation up to a
time-dependent gauge transformation.
\item On solutions, $X_F$ obeys the frame Maxwell system \eqref{eq:RS} with $\bB=\curl\bA$ and $\bE=\curl\mathbf C$.
\item Under $X_A\mapsto X_A+(\nabla\lambda_1;\nabla\lambda_2)$, $X_F$ is unchanged and $\mathcal L$ changes by a divergence.
\item The Hamiltonian is $\tfrac12|\curl X_A|^2=\tfrac12(|\bB|^2+|\bE|^2)$, the energy density.
\item $\mathcal L(e^{\theta G}X_A)=\mathcal L(X_A)$ exactly, and the Noether current of $\delta X_A=GX_A$ has density
$j^0=-\tfrac12X_A\cdot\curl X_A$. For fields that decay or are periodic the conserved charge is $-\mathcal K$, where
\[
\mathcal K=\tfrac12\int X_A\cdot\curl X_A\,d^3x=\tfrac12\int(\bA\cdot\bB+\mathbf C\cdot\bE)\,d^3x
\]
is the optical helicity.
\end{enumerate}
\end{theorem}

This is the duality-symmetric first-order formulation of Schwarz and Sen \cite{SchwarzSen94}, with two vector potentials in the
temporal gauge and with the binding root $G$ in the place of their $SO(2)$ generator; the conserved charge of the duality rotation is,
up to sign, the optical helicity \cite{Calkin65,DeserTeitelboim76,CameronBarnettYao12}.

\begin{example}[The $G$-charge is the helicity]
\label{ex:Gcharge}
For waves $e^{i(k x_3-\omega t)}$ the first-order equation reads $\omega a=G(k\times a)$, a $6\times6$ matrix with eigenvalues $\pm k$, each
twice, and $0$ twice (the longitudinal, pure-gauge pair). The eigenspace $\omega=+k$ is spanned by
\[
\begin{aligned}
a_1&=(1,-i)\otimes(1,i,0), & Ga_1&=ia_1, & s_3a_1&=a_1,\\
a_2&=(1,i)\otimes(1,-i,0), & Ga_2&=-ia_2, & s_3a_2&=-a_2,
\end{aligned}
\]
with $s_3$ acting on each triple. For the real waves $X_A=\Re(a\,e^{ik(x_3-t)})$ the helicity density $\tfrac12X_A\cdot\curl X_A$ and the
energy density are $(k,k^2)$ for $a_1$, $(-k,k^2)$ for $a_2$, and $(0,4k^2\sin^2k(x_3-t))$ for $a_1+a_2$. On positive-frequency modes the
eigenvalue of $G$ is $i$ times the helicity. The ratio of helicity to energy is $\pm1/\omega$, which for a quantised field of energy
$N\hbar\omega$ is $\pm N\hbar$: one unit of $\hbar$ per photon, the spin-one value.
\end{example}

The frame potentials are not components of a Lorentz tensor. Neither the constant boost $\delta X_A=K_3X_A$ nor the boost of the
coordinates is a symmetry of \eqref{eq:SSL}; the boost that is one is the Schwarz--Sen transformation $\delta X_A=t\,\partial_3X_A-x_3\,G\curl
X_A$, whose curl is, identically,
\[
\curl\bigl(t\,\partial_3X_A-x_3\,G\curl X_A\bigr)=t\,\partial_3X_F-x_3\,G\curl X_F-K_3X_F ,
\]
which on solutions is the boost of the field strength along $\hat e_3$ with its orbital part and its spin part $K_3$. Boosts reach the
potentials only through their curl and only on shell \cite{SchwarzSen94}. Duality rotations, for their part, are not Lorentz
transformations: $e^{\theta G}$ preserves the energy density and rotates $\eta(X,X)$ into $\eta(X,GX)$.

\subsection{The bridge between the readings}
\label{sec:bridge}

\begin{proposition}[Half-spinors under the two algebras]
\label{prop:halfspinors}
Under $\lint$ the Casimir on $S^\pm$ is $\tfrac32$, and $S^+$ carries a two-parameter family of invariant antisymmetric forms and no
invariant symmetric one: it is a Weyl module. Under $\lG$ the Casimir on all of $S$ is $3$, the invariant symmetric forms on $S^+$ form a
one-parameter family of Lorentzian forms, and there is no invariant antisymmetric form: $S^+$ is a Minkowski four-vector, and
$\exp(2\pi L_3)=+\id$ on $S$. On $\C^8$ the two algebras give $2\,(\tfrac12,0)\oplus2\,(0,\tfrac12)$ and $2\,(\tfrac12,\tfrac12)$
respectively, and they meet in zero.
\end{proposition}

So under the binding algebra every spinor of $\Cl(3,3)$ has integer spin, in accordance with $\mathrm{Spin}(3,3)\cong SL(4,\R)$, whose
half-spinor modules are $\R^4$ and its dual \cite{Porteous}. This sharpens \cref{prop:emexclusive}: with $\lG$ as the Lorentz algebra
one constituent carries no spin one-half at all.

\begin{theorem}[The Klein correspondence]
\label{thm:klein}
The $\mathrm{Spin}(3,3)$-invariant bilinear forms on $S$ form a two-dimensional space. For one of them, $\beta$, the map
$W(X)(\psi,\chi)=\beta(\psi,X\chi)$ on $S^+$ is antisymmetric and equivariant, and
\[
\Pf W(X)=\eta(X,X).
\]
Hence $V\cong\Lambda^2(S^+)^*$ as $\mathrm{Spin}(3,3)$-modules, and the null cone of $\eta$ is the set of decomposable two-forms.
\end{theorem}

This is the Klein correspondence between the points of the Klein quadric and the lines of projective three-space
\cite{WardWells}. It can be compared with the frame reading. Let $T:\R^{1,3}\to S^+$ be the $\lG$-intertwiner, unique up to scale. Then
\[
T^*W(X)=-F_X+\star F_X=\sqrt2\,F_{e^{3\pi G/4}X}\qquad\text{for all }X,
\]
so the canonical identification of $V$ with two-forms on the half-spinor and the frame reading of the field strength differ by a
duality rotation of an eighth of a turn, $-\pi/4$ modulo $\pi$. In the Klein frame $\eta(X,X)$ is the Pfaffian, the invariant of
$\bE\cdot\bB$ type, and $\eta(X,GX)$ the Hodge pairing; in the frame reading the roles are the other way round. For free fields the two
readings are equivalent, since duality rotations map solutions to solutions; a statement that singles out one invariant, such as which
of them is the Lagrangian, holds in one reading and not in the other. A plane wave $\kappa\wedge\alpha$ built on $S^+$, with $\kappa$ the
lowered null momentum $(1,0,0,1)$ and $\alpha$ the polarisation $(0,1,0,0)$, lands on the doubly null vector
$X=(-\tfrac12,-\tfrac12,0;\,\tfrac12,-\tfrac12,0)$, whose Poynting axis $\bB\times\bE=(0,0,\tfrac12)$ is the spatial momentum; the Klein
planes of $X$ and of $GX$ meet in the null line $[\kappa]$, as two lines of projective space joined by a line of the quadric must.

\begin{figure}[t]
\centering
\begin{tikzpicture}[>=Stealth,font=\small]
\node[draw=tcol,thick,rounded corners=4pt,fill=tcol!6,text width=6.0cm,align=left,inner sep=6pt] (int) at (0,0) {%
\textbf{\textcolor{tcol}{interior reading}}, Lorentz algebra $\lint$\\[3pt]
grade 1: potential $A$, $(\tfrac12,\tfrac12)_0$; surplus pair $(0,0)_{\pm1}$\\
grade 2: field strength $F$, $(1,0)\oplus(0,1)$; $\Omega=\star$\\
spinors: Weyl, $2(\tfrac12,0)\oplus2(0,\tfrac12)$\\
equation: $\nabla F=\mathcal J$ (\cref{thm:nablaF})};
\node[draw=scol,thick,rounded corners=4pt,fill=scol!6,text width=6.0cm,align=left,inner sep=6pt] (bin) at (7.8,0) {%
\textbf{\textcolor{scol}{frame reading}}, Lorentz algebra $\lG$\\[3pt]
grade 1: field strength $X=(\bB;\bE)$; $G=\star$\\
boosts: spin matrices $s_k=ir_k$\\
spinors: four-vectors, $2(\tfrac12,\tfrac12)$\\
equation: $i\partial_tz=(s\cdot p)z$ (\cref{prop:spinmatrices})};
\draw[<->,thick,gray] (int.east) -- node[below,align=center,font=\scriptsize] {not\\ conj.} (bin.west);
\node[draw=gcol,thick,rounded corners=4pt,fill=gcol!6,text width=13.6cm,align=left,inner sep=6pt] (mix) at (3.9,-3.25) {%
\textbf{\textcolor{gcol}{mixed rung $n=2$}}: $\lint$ on constituent 1, $\lG$ on constituent 2. Grade one carries
$V\oplus\Lambda^2V\oplus2\cdot\mathbf1$; the spinor module has a rank-$16$ Dirac block; one grade-three element identifies the
frame field of constituent 2 with $dA$ of constituent 1. Cost: the internal algebra shrinks to one dimension
(\cref{prop:mixedcost}).};
\draw[->,thick,gray] (int.south) -- (int.south |- mix.north);
\draw[->,thick,gray] (bin.south) -- (bin.south |- mix.north);
\node[gray,anchor=south] at (3.9,1.7) {Klein: $V\cong\Lambda^2(S^+)^*$, duality angle $-\pi/4$};
\end{tikzpicture}
\caption{Two places for the photon in the arena (\cref{sec:spin}). Each reading gives a free field with the spin content of the photon;
on one constituent they use non-conjugate Lorentz algebras and do not merge, and at the second rung they can be housed together at
the price of the internal algebra.}
\label{fig:tworoutes}
\end{figure}

\subsection{The tower}

On the tower the two readings multiply in different ways.

\begin{proposition}[Commutants along the tower]
\label{prop:commutantsn}
On $V_n=\R^{3n,3n}$ the commutants have the following dimensions:
\begin{center}
\begin{tabular}{@{}lccccc@{}}
\toprule
commutant of & $n=1$ & $2$ & $3$ & $4$ & $5$\\
\midrule
spectator $\lint$ in $\so{3n,3n}$ & $1$ & $28$ & $91$ & $190$ & $325$\\
diagonal $\lint$ in $\so{3n,3n}$ & $1$ & $7$ & $18$ & $34$ & $55$\\
diagonal $\lG$ in $\so{3n,3n}$ & $0$ & $2$ & $6$ & $12$ & $20$\\
diagonal rotations and $G_n$ in $\so{6n}$ & $1$ & $4$ & $9$ & $16$ & $25$\\
the same inside $\so{3n,3n}$ & $0$ & $1$ & $3$ & $6$ & $10$\\
\bottomrule
\end{tabular}
\end{center}
These are the dimensions of $\so{3n-1,3n-3}$, $\so n\oplus\so{2n}$, $\so{n,\C}$, $\uu n$ and $\so n$.
\end{proposition}

Under the diagonal interior action $V_n=n\,(\tfrac12,\tfrac12)\oplus2n\,(0,0)$, so there are $n$ potentials forming the vector of the copy
algebra $\so n$, and none of them is copy-invariant at grade one for $n\ge2$. A single univectorial potential in the diagonal sector
therefore needs a choice of direction in copy space, and any such choice breaks $\so n$ to $\so{n-1}$. The obstruction belongs to grade
one only.

\begin{theorem}[Copy-invariant four-vectors]
\label{thm:copyparity}
A copy-invariant four-vector in $\Lambda^g(\R^n\otimes\R^{1,3})$ requires $n$ odd and $g$ odd with $n\le g\le3n$, and for every odd $n$
there is one at $g=n$,
\[
Q^{(n)}_b=\epsilon_{k_1\cdots k_n}\,\eta^{a_1a_2}\cdots\eta^{a_{n-2}a_{n-1}}\,\gamma^{(k_1)}_{a_1}\wedge\cdots\wedge\gamma^{(k_{n-1})}_{a_{n-1}}\wedge\gamma^{(k_n)}_b .
\]
At $n=3$ there is exactly one at each of the grades $3,5,7,9$, and at $n=2$ there is none at any grade.
\end{theorem}

\begin{proof}
An invariant element of $\Lambda^g(\R^n\otimes\R^{1,3})\otimes\R^{1,3}$ is an invariant tensor of $SO(n)\times SO^+(1,3)$ in
$(\R^n)^{\otimes g}\otimes(\R^{1,3})^{\otimes(g+1)}$. By the first fundamental theorem for the special orthogonal groups \cite{Weyl39},
such tensors are spanned by products of metrics and at most one Levi-Civita symbol. On the Lorentz side every factor has an even number
of indices, so $g+1$ is even. On the copy side $g=2d+ne$ with $e\in\{0,1\}$; $g$ odd forces $e=1$ and $n$ odd, so $g\ge n$. The Hodge
star of $\R^n\otimes\R^{1,3}$ is an equivariant isomorphism $\Lambda^g\cong\Lambda^{4n-g}$, which gives $g\le3n$. The displayed
$Q^{(n)}_b$ is invariant because $\epsilon$ is, and nonzero because the terms of $Q^{(n)}_0$ with all Lorentz indices equal to $0$
add up to $n!\,\gamma^{(1)}_0\wedge\cdots\wedge\gamma^{(n)}_0$.
\end{proof}

At $n=3$ the invariant four-vector of grade three is, up to a factor, the copy-singlet cubic $Q_\mu$ that builds the gamma matrices of
the Dirac block in \cref{thm:blockgamma}; the parity theorem explains why such an object exists at three constituents and could not
exist at two. Under the diagonal $\lG$, finally, $V_n=n\bigl[(1,0)\oplus(0,1)\bigr]$: $n$ field strengths mixed by the complex orthogonal
flavour algebra $\so{n,\C}$. The energy-preserving duality group of $n$ free Maxwell fields is $U(n)$, the maximal compact subgroup of
the Gaillard--Zumino group $Sp(2n,\R)$ \cite{GaillardZumino81}; its Lie algebra meets $\so{3n,3n}$ only in $\so n$, and its centre,
generated by $G_n$ and carrying the total helicity, lies outside $\so{3n,3n}$, as $G$ does at $n=1$. A single photon would be one
direction in this flavour space, and nothing selects one.

\subsection{The mixed rung}
\label{sec:mixed}

\Cref{prop:emexclusive} forbids reading one six-dimensional space in two ways. At $n\ge2$ there is more than one constituent, and the
obvious repair is to let one constituent carry the interior reading and another the frame reading, under one Lorentz action.

\begin{proposition}[The mixed embedding]
\label{prop:mixed}
In $\Cl(6,6)$ let $\mathfrak l_{\rm mix}$ be spanned by $S_k\otimes1+1\otimes L_k$ and $T_k\otimes1+1\otimes K_k$: the interior algebra on
constituent $1$ and $\lG$ on constituent $2$, paired generator by generator. It satisfies the $\so{1,3}$ brackets, and
\[
\R^{6,6}=\underbrace{V}_{C=3}\oplus\underbrace{\Lambda^2V}_{C=4}\oplus2\cdot\mathbf1 ,
\]
so one Lorentz action carries both a four-vector and a field strength at grade one. On $S_2=\R^{64}$ the Casimir satisfies
$(C_{\rm mix}-\tfrac32)(C_{\rm mix}-\tfrac{11}2)=0$, with the Dirac value $\tfrac32$ on sixteen dimensions and $\tfrac{11}2$ on $48$.
\end{proposition}

\begin{proposition}[The cost]
\label{prop:mixedcost}
The centraliser of $\mathfrak l_{\rm mix}$ in $\so{6,6}$ is one-dimensional, spanned by the rotation of the surplus pair of
constituent $1$. At the same rung:
\begin{center}
\begin{tabular}{@{}lcl@{}}
\toprule
Lorentz action at $n=2$ & centraliser & content of $\R^{6,6}$\\
\midrule
interior on one constituent & $28$ & $V\oplus8\cdot\mathbf1$\\
$\lG$ on one constituent & $15$ & $\Lambda^2V\oplus6\cdot\mathbf1$\\
mixed, across both & $1$ & $V\oplus\Lambda^2V\oplus2\cdot\mathbf1$\\
\bottomrule
\end{tabular}
\end{center}
The mixed embedding buys the field-strength module at the price of the whole internal algebra: $\so{5,3}$ collapses to one $\uu1$.
\end{proposition}

The count is for this pairing of the two Lorentz bases at $n=2$; another pairing or a higher rung could do better and has not been
examined. Within its scope, however, the mixed rung is an honest place to test whether the two readings can be coupled, and it turns out
that they can.

\begin{theorem}[The Dirac block of the mixed rung]
\label{thm:mixeddirac}
Let $P=\tfrac18(11\,\id-2C_{\rm mix})$ and $\hat\gamma_\mu=2P\gamma^{(1)}_\mu P$, with $J_1=e_2e_3\otimes1$.
\begin{enumerate}[label=(\roman*),leftmargin=2.2em]
\item $P$ is the orthogonal projector of rank $16$ onto the block with the Dirac value, and commutes with the mixed generators.
\item $\{\hat\gamma_\mu,\hat\gamma_\nu\}=2\eta_{\mu\nu}P$, $[\hat\gamma_\mu,J_1]=0$, $J_1^2=-\id$, and $P\mathcal S_kP=\tfrac12\hat\gamma_i\hat\gamma_j$ (cyclic),
$P\mathcal T_kP=\tfrac12\hat\gamma_0\hat\gamma_k$: the compressed gammas generate the mixed Lorentz action on the block.
\item With $J_1$ as the imaginary unit the block is two complex Dirac modules, labelled by the half-spinors of constituent $2$, with
flavour algebra $\uu2$; $\mathcal H(p)=\hat\gamma^0\hat\gamma^ap_a+m\hat\gamma^0$ is symmetric and $\mathcal H(p)^2=(|p|^2+m^2)P$.
\item The block is the $\gamma$-trace part $\sum_a\gamma_a\chi\otimes\tau e_a$ of the vector-spinors (Dirac spinor of constituent $1$)
$\otimes$(four-vector of constituent $2$), and the rank-$48$ block is their $\gamma$-traceless $(1,\tfrac12)\oplus(\tfrac12,1)$ part.
\end{enumerate}
\end{theorem}

Item (iv) explains the factor $2$ in $\hat\gamma_\mu$, since $P\gamma^{(1)}_\mu$ acts on the $\gamma$-trace part as $-\tfrac12\gamma_\mu$, and it
identifies the decomposition with the classical splitting of a vector-spinor into its trace and its traceless part
\cite{RaritaSchwinger41}.

A covariant identification of the field strength of constituent $1$ with the frame field of constituent $2$ is an invariant element
of $\Lambda^2(\text{four-plane}_1)\otimes V_2$, that is, a grade-three element of the algebra commuting with the mixed generators. The
invariant space is two-dimensional, and one invariant is
\[
\Theta_F=\gamma^{(1)}_2\gamma^{(1)}_3e'_1+\gamma^{(1)}_3\gamma^{(1)}_1e'_2+\gamma^{(1)}_1\gamma^{(1)}_2e'_3
-\gamma^{(1)}_0\gamma^{(1)}_1e'_4-\gamma^{(1)}_0\gamma^{(1)}_2e'_5-\gamma^{(1)}_0\gamma^{(1)}_3e'_6 ,
\]
with $e'_a$ the generators of constituent $2$; it pairs the spatial bivectors of constituent $1$ with the magnetic directions of
constituent $2$ and the time--space bivectors with the electric ones. The other invariant is $\Theta_F$ followed by $G$, so up to scale
the identifications form the family $X=e^{\theta G}X_F$.

\begin{theorem}[A coupled model]
\label{thm:coupled}
On the block, in flavour coordinates $\R^8\otimes\R^2$ with $\hat\gamma_\mu=-\gamma_\mu\otimes1$, $J_1=J\otimes1$ and $U^\mu=\hat\gamma^0\hat\gamma^\mu$, take
for commuting real fields
\begin{equation}
\label{eq:coupledL}
\mathcal L=-\tfrac14F_{\mu\nu}F^{\mu\nu}+\psi^{\!\top}J_1U^\mu D_\mu\psi-m\,\psi^{\!\top}\hat\gamma^0\psi,\qquad D_\mu=\partial_\mu+qA_\mu J_1 ,
\end{equation}
the real form of $-\tfrac14F^2+\bar\psi(i\gamma^\mu D_\mu-m)\psi$ with $J_1$ in the place of $i$. Then the field equations are Maxwell's
equations $\partial_\mu F^{\mu\nu}=q\,j^\nu$ with $j^\nu=\psi^{\!\top}U^\nu\psi$ and $j^0=|\psi|^2$, and the Dirac equation
$J_1\hat\gamma^\mu D_\mu\psi=m\psi$; the identity $\partial_\mu j^\mu+\psi^{\!\top}J_1\mathcal E_\psi=0$ holds with $\mathcal E_\psi$ the Dirac
expression; the Lagrangian is exactly invariant under $\psi\mapsto e^{-q\lambda J_1}\psi$, $A\mapsto A+\partial\lambda$; and its flavour
symmetry is the $\uu2$ spanned by $J\otimes1$, $1\otimes\epsilon$, $J\otimes\sigma_x$, $J\otimes\sigma_z$. The potential carries two
polarisations of helicity $\pm1$ and no massive mode, and each momentum carries two flavours, two helicities and two signs of frequency.
\end{theorem}

In the frame reading the source is electric. For the frame field of an arbitrary potential one has, identically,
$\partial_tz+i\curl z+i\boldsymbol{\mathcal J}=0$ and $\nabla\cdot z=i\rho$ with $(\rho,\boldsymbol{\mathcal J})=\partial_\mu F^{\mu\nu}$, and the
rotated field $e^{\theta G}X_F$ satisfies the same system with electric source $\cos\theta$ times the current and magnetic source $\sin\theta$
times it. Since the frame field takes no part in the dynamics of \eqref{eq:coupledL}, which member of the family is called the frame
field is a choice; requiring it to carry electric sources only fixes $\theta\equiv0$ modulo $\pi$, as in the textbook treatment of duality
with sources, where the angle is fixed by the convention that all charges are electric \cite{Jackson}. The duality rotation, a symmetry
of the free system and of \eqref{eq:SSL}, is not a symmetry once electric sources are present.

The coupled model is classical two-flavour spinor electrodynamics with equal charges, written in the algebra of the second rung. It
is a construction, not a selection, and it has a definite limitation that we record now and use later: the second-rung vacuum of
\cref{sec:photon} is defined for the spectator action, and under the mixed action it is preserved only by the rotations about one axis,
so the coupled model lies outside that vacuum sector.

At even rungs, finally, the diagonal interior action places integer-spin modules inside the spinor module itself: at $n=2$ the Casimir
takes the values $0,3,4$ on $8,32,24$ dimensions, so $S_2$ contains scalars of copy charge $\pm2$, four-vectors of copy charge $\pm1$ and
copy-neutral field strengths $(1,0)\oplus(0,1)$, and the copy and Lorentz Casimirs add up to $4$, the $n=2$ case of \eqref{eq:casid}. These
are modules of one spinor space, not bound states, and we do not propose them as the photon.
